% ============================
%   6. Current state
% ============================
\section{Current state of the project}
\label{sec:current_state}

\subsection{Functional milestones reached}

\begin{enumerate}
    \item \textbf{Backend skeleton.} FastAPI application with async
        SQLAlchemy, JWT auth, OpenAPI documentation auto-served at
        \code{/docs}.
    \item \textbf{Signature engine.} YAML seed file, DB-backed
        editable rules, hot reload from the API, six bundled
        categories (Recon, BruteForce, Injection, RemoteAccess,
        DenialOfService, DataExfiltration).
    \item \textbf{Capture layer.} Both LiveSniffer (Scapy) and the
        synthetic TrafficSimulator are functional; the simulator
        currently drives 99\% of demonstration usage because raw
        capture requires elevated privileges.
    \item \textbf{ML detector.} Random Forest and Isolation Forest can
        be trained from the UI on a synthetic feature set; metrics,
        feature importance and the active model are surfaced on the
        ML Models page.
    \item \textbf{Forensics.} Per-alert JSONL ring captures are
        produced, listed and downloadable; the Alert Detail page
        renders the packet timeline for replay.
    \item \textbf{Frontend.} Login, Dashboard, Alerts, Alert Detail,
        Rules, Forensics, ML Models and Settings pages are
        feature-complete for the current milestone. Live updates via
        WebSocket cover the dashboard chart and the alert stream.
    \item \textbf{Deployment.} The one-click launcher works on
        Windows, macOS and Linux; \code{run.bat} is a double-click
        entry point and the auto-installer creates the venv and
        installs both Python and Node dependencies on first launch.
\end{enumerate}

\subsection{Quantitative observations}

The figures below are not formal benchmarks; they were collected
during regular development on a laptop (Windows 11, Intel i7,
16~GB RAM, Python 3.12, Node 20) and are reported here only as
sanity checks for the upcoming evaluation chapter.

\begin{table}[h]
\centering
\small
\begin{tabularx}{\textwidth}{l X}
\toprule
\textbf{Observation} & \textbf{Value} \\
\midrule
Steady-state simulator throughput & 24--40 packets/s \\
Burst peak (during an attack injection) & 90--120 packets/s \\
Median signature-engine evaluation time per packet & $< 1$\,ms \\
Median ML score (Random Forest) per packet & $\approx 0.4$\,ms \\
Random Forest training time (synthetic, 4000 samples) & $\approx 1.0$\,s \\
Isolation Forest training time (same dataset) & $\approx 0.7$\,s \\
Random Forest accuracy on held-out 25\% split & 0.99 \\
Isolation Forest accuracy on held-out 25\% split & 0.83 \\
End-to-end alert latency (packet $\rightarrow$ UI) & $< 250$\,ms \\
Backend cold start (first request served) & 1.5--3\,s \\
Frontend cold start (Vite dev server) & 2--5\,s \\
\bottomrule
\end{tabularx}
\caption{Indicative measurements during day-to-day usage.}
\end{table}

The unsupervised Isolation Forest unsurprisingly trails the supervised
Random Forest on this synthetic dataset where labels are available; its
purpose is to provide a deployment path when labelled attack data is
not available, not to win an accuracy contest.

\subsection{Known limitations}

\begin{itemize}
    \item Live packet capture has been tested only on Linux; Windows
        and macOS require manual installation of WinPcap/Npcap or
        equivalent, which the launcher does not yet automate.
    \item The ML feature vector is currently per-packet rather than
        per-flow; flow-aware aggregation is on the roadmap for the
        final thesis.
    \item No clustering or correlation engine yet; alerts are surfaced
        individually. A future ``alert grouping by source IP'' view is
        sketched.
    \item Database migrations rely on \code{create\_all}. An Alembic
        suite is planned before any production-style deployment.
    \item The simulator's attack patterns are deterministic and
        comparatively easy for the ML detector to flag. Evaluation on
        the CICIDS2017~\cite{sharafaldin2018cicids} dataset is planned
        for the experimental chapter.
\end{itemize}

\subsection{Repository state}

At the time of writing, the repository contains approximately:
\begin{itemize}
    \item 1{,}900 lines of backend Python (engine + API + models +
        schemas + tests scaffolding).
    \item 2{,}600 lines of frontend JavaScript/JSX across eight pages
        and the shared component library.
    \item 480 lines of YAML/JSON for the default rule base and
        configuration.
    \item A single-pass-compilable LaTeX folder (this document) with
        bibliography and section split.
\end{itemize}
The default credentials in the seeded admin user are
\code{admin}/\code{admin}; the password can now be rotated from the
Settings page after first login.
